\ifdefined\gkver\else\def\gkver{student}\fi
\documentclass[\gkver]{gaokaozhenti}
\begin{document}

\chapter{1989年上海}

\begin{enumerate}
\item
%% number: 1
%% paperName: 1989年上海高考第 1 题 4 分
%% typeId: 1
%% type: 单选题
%% chapter: 光学
%% point: 光的电磁说
%% method: 无
%% score: 4
%% degree: 850
%% duplicateId: 0
%% body:
下列各组电磁波，按波长由长到短正确排列的是\xzanswer{B}
\fourchoices[answer=B]
{$\gamma$ 射线、红外线、紫外线、可见光}
{红外线、可见光、紫外线、$\gamma$ 射线}
{可见光、红外线、紫外线、$\gamma$ 射线}
{紫外线、可见光、红外线、$\gamma$ 射线}

%% answer: B
%% memo:
\memoanswer{
本题原卷及材料中未提供详解，待补充。
}

\item
%% number: 2
%% paperName: 1989年上海高考第 2 题 4 分
%% typeId: 1
%% type: 单选题
%% chapter: 固体液体的性质
%% point: 固体的基本性质
%% method: 无
%% score: 4
%% degree: 850
%% duplicateId: 0
%% body:
下列各组物质，全部都是晶体的是\xzanswer{C}
\fourchoices[answer=C]
{石英、雪花、沥青}
{食盐、橡胶、沥青}
{食盐、雪花、石英}
{雪花、橡胶、石英}

%% answer: C
%% memo:
\memoanswer{
本题原卷及材料中未提供详解，待补充。
}

\item
%% number: 3
%% paperName: 1989年上海高考第 3 题 4 分
%% typeId: 1
%% type: 单选题
%% chapter: 机械振动
%% point: 简谐振动
%% method: 无
%% score: 4
%% degree: 650
%% duplicateId: 0
%% body:
弹簧振子沿直线作简谐振动。当振子连续两次经过平衡位置时，振子的\xzanswer{A}
\fourchoices[answer=A]
{加速度相同，动能相同}
{动能相同，动量相同}
{加速度相同，速度相同}
{动量相同，速度相同}

%% answer: A
%% memo:
\memoanswer{
本题原卷及材料中未提供详解，待补充。
}

\item
%% number: 4
%% paperName: 1989年上海高考第 4 题 4 分
%% typeId: 1
%% type: 单选题
%% chapter: 直线运动
%% point: 竖直上下抛运动
%% method: 无
%% score: 4
%% degree: 650
%% duplicateId: 0
%% body:
升降机以加速度 $a$ 竖直向上作匀加速运动。升降机内的天花板上有一只螺帽突然松动，脱离天花板。这时螺帽相对于地的加速度是（$g$ 为重力加速度）\xzanswer{D}
\fourchoices[answer=D]
{$g-a$}
{$g+a$}
{$a$}
{$g$}

%% answer: D
%% memo:
\memoanswer{
本题原卷及材料中未提供详解，待补充。
}

\item
%% number: 5
%% paperName: 1989年上海高考第 5 题 4 分
%% typeId: 1
%% type: 单选题
%% chapter: 交流电
%% point: 变压器
%% method: 无
%% score: 4
%% degree: 650
%% duplicateId: 0
%% body:
右图为一与电源相接的理想降压变压器。原线圈中电流为 $I_{1}$，副线圈中电流为 $I_{2}$。当副线圈中的负载电阻 $R$ 变小时\xzanswer{C}
\onepicture[w=7cm]{\includesvg{figs/05}}
\fourchoices[answer=C]
{$I_{2}$ 变小，$I_{1}$ 变小}
{$I_{2}$ 变小，$I_{1}$ 增大}
{$I_{2}$ 增大，$I_{1}$ 增大}
{$I_{2}$ 增大，$I_{1}$ 变小}

%% answer: C
%% memo:
\memoanswer{
本题原卷及材料中未提供详解，待补充。
}

\gknewpage

\item
%% number: 6
%% paperName: 1989年上海高考第 6 题 4 分
%% typeId: 1
%% type: 单选题
%% chapter: 气体性质
%% point: 气体图象
%% method: 无
%% score: 4
%% degree: 650
%% duplicateId: 0
%% body:
一定质量的理想气体，经历了 A$\to$B$\to$C$\to$D 的状态变化过程，如右边的 $p-V$ 图所示，其中 BC 段是以 $p$ 轴和 $V$ 轴为渐近线的双曲线。在 $p$（压强）$-T$（温度）图上，上述过程可以对应为\xzanswer{B}
\onepicture[w=4.6cm]{\includesvg{figs/06a}}
\fourchoices[answer=B, ispicture=true, h=2.2cm]
{\includesvg{figs/06b}}
{\includesvg{figs/06c}}
{\includesvg{figs/06d}}
{\includesvg{figs/06e}}

%% answer: B
%% memo:
\memoanswer{
本题原卷及材料中未提供详解，待补充。
}

\item
%% number: 7
%% paperName: 1989年上海高考第 7 题 4 分
%% typeId: 1
%% type: 单选题
%% chapter: 电磁感应
%% point: 楞次定律推广
%% method: 无
%% score: 4
%% degree: 650
%% duplicateId: 0
%% body:
图示一通有电流 $I$ 的直导线和一矩形导线框平行放置在同一平面上，当线框向哪个方向运动时，才会受到向右的合力\xzanswer{D}
\onepicture[h=6cm]{\includesvg{figs/07}}
\fourchoices[answer=D]
{向上}
{向下}
{向右}
{向左}

%% answer: D
%% memo:
\memoanswer{
本题原卷及材料中未提供详解，待补充。
}

\item
%% number: 8
%% paperName: 1989年上海高考第 8 题 4 分
%% typeId: 1
%% type: 单选题
%% chapter: 曲线运动
%% point: 万有引力定律
%% method: 无
%% score: 4
%% degree: 650
%% duplicateId: 0
%% body:
地球表面的平均重力加速度为 $g$，地球半径为 $R$，万有引力恒量为 $G$。可以用下式来估计地球的平均密度\xzanswer{A}
\fourchoices[answer=A]
{$\frac{3g}{4\pi RG}$}
{$\frac{3g}{4\pi R^{2}G}$}
{$\frac{g}{RG}$}
{$\frac{g}{R^{2}G}$}

%% answer: A
%% memo:
\memoanswer{
本题原卷及材料中未提供详解，待补充。
}

\item
%% number: 9
%% paperName: 1989年上海高考第 9 题 5 分
%% typeId: 2
%% type: 多选题
%% chapter: 气体性质
%% point: 物体的内能
%% method: 无
%% score: 5
%% degree: 650
%% duplicateId: 0
%% body:
封闭在体积一定的容器内的理想气体，当温度升高时，下列说法中正确的是\xzanswer{BC}
\fourchoices[answer=BC]
{气体分子的密度增加}
{气体分子的平均动能增加}
{气体分子的平均速率增加}
{气体分子的势能增加}

%% answer: BC
%% memo:
\memoanswer{
本题原卷及材料中未提供详解，待补充。
}

\gknewpage

\item
%% number: 10
%% paperName: 1989年上海高考第 10 题 5 分
%% typeId: 2
%% type: 多选题
%% chapter: 交流电
%% point: 交流电
%% method: 无
%% score: 5
%% degree: 650
%% duplicateId: 0
%% body:
一个按正弦规律变化的交流电流的图象如下图所示，根据图象可以知道\xzanswer{BD}
\onepicture[w=6.5cm]{\includesvg{figs/10}}
\fourchoices[answer=BD]
{该交流电流的频率是 $0.02\UHz$}
{该交流电流的有效值是 $14.1\UA$}
{该交流电流的瞬时值表示式是 $i=20\sin0.02t\UA$}
{在 $t=\frac{T}{8}$（$T$ 是周期）时刻，该电流的大小与其有效值相等}

%% answer: BD
%% memo:
\memoanswer{
本题原卷及材料中未提供详解，待补充。
}

\item
%% number: 11
%% paperName: 1989年上海高考第 11 题 5 分
%% typeId: 1
%% type: 单选题
%% chapter: 原子和原子核
%% point: 半衰期
%% method: 无
%% score: 5
%% degree: 650
%% duplicateId: 0
%% body:
用下述方法可以减缓放射性元素的衰变。\xzanswer{D}
\fourchoices[answer=D]
{把该元素放置在低温处}
{把该元素密封在很厚的铅盒子里}
{把该元素同其他的稳定元素结合成化合物}
{上述各种办法无法减缓放射性元素的衰变}

%% answer: D
%% memo:
\memoanswer{
各种放射性元素的衰变的快慢程度是不同的，这是它们的一个特性，基本上不随外界条件（温度、压强等）的变化而变化，并与元素所处的状态（游离态还是化合态）无关。
}

\item
%% number: 12
%% paperName: 1989年上海高考第 12 题 5 分
%% typeId: 2
%% type: 多选题
%% chapter: 光学
%% point: 透镜
%% method: 无
%% score: 5
%% degree: 650
%% duplicateId: 0
%% body:
凸透镜的焦距为 $f$。一个在透镜光轴上的物体，从距离透镜 $3f$ 处，沿光轴逐渐移动到距透镜 $1.5f$ 处，在此过程中\xzanswer{AC}
\fourchoices[answer=AC]
{像不断增大}
{像和物之间的距离不断增大}
{像和焦点的距离不断增大}
{像和透镜的距离不断减小}

%% answer: AC
%% memo:
\memoanswer{
本题原卷及材料中未提供详解，待补充。
}

\item
%% number: 13
%% paperName: 1989年上海高考第 13 题 5 分
%% typeId: 2
%% type: 多选题
%% chapter: 电场
%% point: 带电粒子的平衡
%% method: 无
%% score: 5
%% degree: 450
%% duplicateId: 0
%% body:
平行板电容器的两极板 A、B 接于电池两极。一带正电小球悬挂在电容器内部。闭合电键 K，电容器充电，这时悬线偏离竖直方向的夹角为 $\theta$，如右图所示\xzanswer{AD}
\onepicture[h=7cm]{\includesvg{figs/13}}
\fourchoices[answer=AD]
{保持电键 K 闭合，带正电的 A 板向 B 板靠近，则 $\theta$ 增大}
{保持电键 K 闭合，带正电的 A 板向 B 板靠近，则 $\theta$ 不变}
{电键 K 断开，带正电的 A 板向 B 板靠近，则 $\theta$ 增大}
{电键 K 断开，带正电的 A 板向 B 板靠近，则 $\theta$ 不变}

%% answer: AD
%% memo:
\memoanswer{
本题原卷及材料中未提供详解，待补充。
}

\item
%% number: 14
%% paperName: 1989年上海高考第 14 题 4 分
%% typeId: 1
%% type: 单选题
%% chapter: 气体性质
%% point: 能量的转化及其方向性
%% method: 无
%% score: 4
%% degree: 850
%% duplicateId: 0
%% body:
能量既不能凭空产生，也不能凭空消灭，它只能\tkanswer{从一种形式转化为别的形式}，或者\tkanswer{从一个物体转移到别的物体}，能量的总和保持不变，这就是能的转化和守恒定律。

%% answer: 
%% memo:
\memoanswer{
本题原卷及材料中未提供详解，待补充。
}

\item
%% number: 15
%% paperName: 1989年上海高考第 15 题 4 分
%% typeId: 3
%% type: 填空题
%% chapter: 气体性质
%% point: 盖吕萨克定律
%% method: 无
%% score: 4
%% degree: 650
%% duplicateId: 0
%% body:
用长为 $h_{0}=50.0\Umm$ 的一段水银柱，把空气封闭在一端开口向上的粗细均匀的玻璃管内，气柱高度为 $h_{1}=27.3\Umm$，室温为 $273\UK$，大气压强 $p_{0}=760\UmmHg$。这时，封闭在玻璃管内的空气柱压强 $p=$\tkanswer{$810\UmmHg$}，当室温升高 $20\UdC$ 时，封闭在管内的空气柱高度 $h=$\tkanswer{$29.3\Umm$}。

%% answer: 810，29.3
%% memo:
\memoanswer{
本题原卷及材料中未提供详解，待补充。
}

\item
%% number: 16
%% paperName: 1989年上海高考第 16 题 4 分
%% typeId: 3
%% type: 填空题
%% chapter: 力物体的平衡
%% point: 三力平衡
%% method: 临界
%% score: 4
%% degree: 650
%% duplicateId: 0
%% body:
如图，在一细绳 C 点系住一重物 P，细绳两端 A、B 分别固定在墙面上。使得 AC 保持水平，BC 与水平方向成 $30^{\circ}$ 角。已知细绳最大只能承受 $200\UN$ 的拉力，那么 C 点悬挂物的重量最多为\tkanswer{$100\UN$}，这时细绳的\tkanswer{BC}段即将断裂。
\onepicture[w=7cm, align=right]{\includesvg{figs/16}}

%% answer: 100，BC
%% memo:
\memoanswer{
本题原卷及材料中未提供详解，待补充。
}

\item
%% number: 17
%% paperName: 1989年上海高考第 17 题 4 分
%% typeId: 3
%% type: 填空题
%% chapter: 光学
%% point: 透镜
%% method: 无
%% score: 4
%% degree: 650
%% duplicateId: 0
%% body:
一摄影者使用焦距是 $5.00\Ucm$ 的照相机拍摄距离 $1.05\Um$ 处的一个台灯，在冲洗好的底片上量得此台灯的高度是 $2.00\Ucm$，此台灯的实际高度是\tkanswer{$40\Ucm$}。

%% answer: 40 厘米
%% memo:
\memoanswer{
本题原卷及材料中未提供详解，待补充。
}

\item
%% number: 18
%% paperName: 1989年上海高考第 18 题 4 分
%% typeId: 3
%% type: 填空题
%% chapter: 直线运动
%% point: 平均速度和瞬时速度
%% method: 无
%% score: 4
%% degree: 650
%% duplicateId: 0
%% body:
一物体作同向直线运动，前一半时间以 $9.0\Ums$ 的速度作匀速运动，后一半时间以 $6.0\Ums$ 的速度作匀速运动，则物体的平均速度是\tkanswer{$7.5\Ums$}。另一物体也作同向直线运动，前一半路程以 $3.0\Ums$ 的速度作匀速运动；后一半路程以 $7.0\Ums$ 的速度作匀速运动，则物体的平均速度是\tkanswer{$4.2\Ums$}。

%% answer: 7.5，4.2
%% memo:
\memoanswer{
本题原卷及材料中未提供详解，待补充。
}

\item
%% number: 19
%% paperName: 1989年上海高考第 19 题 4 分
%% typeId: 3
%% type: 填空题
%% chapter: 电场
%% point: 带电粒子的平衡
%% method: 无
%% score: 4
%% degree: 650
%% duplicateId: 0
%% body:
如图所示，电量分别为 $Q$ 和 $q$ 的两个点电荷，放置在 A 点和 B 点，已知 AB $=$ BC $=L$，DB $=\frac{L}{4}$。现在把一点电荷 $q_{0}$ 依次放到 C 点和 D 点，为了使两电荷对 $q_{0}$ 的两个电场力大小相等，$q_{0}$ 在 C 点时，$\left|\frac{Q}{q}\right|=$\tkanswer{4:1}；$q_{0}$ 在 D 点时，$\left|\frac{Q}{q}\right|=$\tkanswer{9:1}。
\onepicture[w=8cm, align=right]{\includesvg{figs/19}}

%% answer: 4:1，9:1
%% memo:
\memoanswer{
本题原卷及材料中未提供详解，待补充。
}

\item
%% number: 20
%% paperName: 1989年上海高考第 20 题 4 分
%% typeId: 3
%% type: 填空题
%% chapter: 光学
%% point: 光电效应
%% method: 无
%% score: 4
%% degree: 650
%% duplicateId: 0
%% body:
用同一束单色光，在同一条件下，先后照射锌片和银片，都能产生光电效应。对于这两个过程，下列括号所列的四个物理量中，一定相同的是\tkanswer{①}，可能相同的是\tkanswer{③}，一定不同的是\tkanswer{②④}。（只须填各量前的编号）

[①照射光子的能量；②光电子的逸出功；③光电子的动能；④光电子的最大初速度。]

%% answer: ①，③，②④
%% memo:
\memoanswer{
本题原卷及材料中未提供详解，待补充。
}

\item
%% number: 21
%% paperName: 1989年上海高考第 21 题 4 分
%% typeId: 3
%% type: 填空题
%% chapter: 运动和力
%% point: 动量守恒定律
%% method: 无
%% score: 4
%% degree: 450
%% duplicateId: 0
%% body:
长为 $l$ 的轻绳，一端用轻环套在水平光滑的横杆上（轻绳与轻环的质量都忽略不计），另一端连接一质量为 $m$ 的小球。开始时，将系球的绳子绷紧并转到与横杆平行位置，然后轻轻放手。当绳子与横杆成 $\theta$ 角时（见右图），小球速度在水平方向的分量大小是\tkanswer{0}，竖直方向的分量大小是\tkanswer{$\sqrt{2gl\sin\theta}$}。
\onepicture[w=6cm, align=right]{\includesvg{figs/21}}

%% answer: 0，$\sqrt {2gl\sin \theta }$
%% memo:
\memoanswer{
本题原卷及材料中未提供详解，待补充。
}

\item
%% number: 22
%% paperName: 1989年上海高考第 22 题 6 分
%% typeId: 4
%% type: 实验题
%% chapter: 运动和力
%% point: 验证牛顿第二定律
%% method: 无
%% score: 6
%% degree: 650
%% duplicateId: 0
%% body:
在做“验证牛顿第二定律”的实验时，为了平衡摩擦力，需要在长木板的下面垫一木块（木块垫在长木板的不带定滑轮的一端）。反复移动木块的位置，直到测出小车所拖纸带上的各个相邻记数点之间的距离都\tkanswer{相等}时为止。这时，小车在斜面上所做的是\tkanswer{匀速}运动，小车拖着纸带运动时受到的摩擦阻力恰好与小车的\tkanswer{重力在斜面方向上的分力}平衡。

%% answer: 相等，匀速，重力在斜面方向上的分力。
%% memo:
\memoanswer{
本题原卷及材料中未提供详解，待补充。
}

\item
%% number: 23
%% paperName: 1989年上海高考第 23 题 6 分
%% typeId: 4
%% type: 实验题
%% chapter: 电路
%% point: 电路实验
%% method: 无
%% score: 6
%% degree: 650
%% duplicateId: 0
%% body:
在“把电流表改装为伏特表”的实验中，利用右图所示的电路来测定电流表的内阻 $r_{g}$。为此，应进行下列步骤的操作：

①使电键 $K_{1}$ 闭合，$K_{2}$\tkanswer{断开}，调节\tkanswer{$R$}，电流表指针偏转到满刻度处。

②电键 $K_{1}$ 仍闭合，$K_{2}$\tkanswer{闭合}，固定\tkanswer{$R$}，调节\tkanswer{$R'$}，使电流表指针偏转到正好是满刻度的一半处。

③断开电键 $K_{1}$ 读出电阻箱 $R'$ 的阻值。在 $R$ 与 $R'$ 满足\tkanswer{$R\gg R'$}时，电流表的内阻 $r_{g}$ 可以认为等于 $R'$。
\onepicture[w=6.5cm, align=right]{\includesvg{figs/23}}

%% answer: （1）断开，R；（2）闭合，R，R'；（3）R ≫ R'
%% memo:
\memoanswer{
本题原卷及材料中未提供详解，待补充。
}

\item
%% number: 24
%% paperName: 1989年上海高考第 24 题 4 分
%% typeId: 4
%% type: 实验题
%% chapter: 电磁感应
%% point: 自感与互感，涡流
%% method: 无
%% score: 4
%% degree: 650
%% duplicateId: 0
%% body:
右图是演示自感现象的实验电路图，$L$ 是电感线圈，$A_{1}$、$A_{2}$ 是规格相同的灯泡，$R$ 的阻值与 $L$ 的电阻值相同。当开关 K 由断开到合上时，观察到自感现象是\tkanswer{灯泡 $A_{2}$ 立刻正常发光，灯泡 $A_{1}$ 逐渐亮起来，（或 $A_{2}$ 先亮，$A_{1}$ 后亮）}，最后达到同样亮。
\onepicture[w=6cm, align=right]{\includesvg{figs/24}}

%% answer: 灯泡 A_2 立刻正常发光，灯泡 A_1 逐渐亮起来，（或 A_2 先亮，A_1 后亮）
%% memo:
\memoanswer{
本题原卷及材料中未提供详解，待补充。
}

\item
%% number: 25
%% paperName: 1989年上海高考第 25 题 6 分
%% typeId: 4
%% type: 实验题
%% chapter: 电路
%% point: 多用表的使用
%% method: 无
%% score: 6
%% degree: 650
%% duplicateId: 0
%% body:
用万用电表的欧姆档检测晶体二极管的好坏。

①将表笔 c、d 分别与二级管 a、b 端接触，其结果如图（A）所示。

②对调表笔所接触的端点，再进行上述检测，若结果如图（B）所示，则可以判定此二极管是\tkanswer{好}的，若结果如图（C）所示，则可以判定该二极管是\tkanswer{坏}的。

③可以判定，完好的二级管的 a 端是\tkanswer{负}极，b 端是\tkanswer{正}极。
\onepicture[w=10cm, align=right]{\includesvg{figs/25}}

%% answer: ②好；坏。③负；正。
%% memo:
\memoanswer{
本题原卷及材料中未提供详解，待补充。
}

\item
%% number: 26
%% paperName: 1989年上海高考第 26 题 10 分
%% typeId: 6
%% type: 计算题
%% chapter: 运动和力
%% point: 牛顿第二定律
%% method: 无
%% score: 10
%% degree: 850
%% duplicateId: 0
%% body:
质量 $m=5.0\Ukg$ 的物体，置于倾角 $\alpha=30^{\circ}$ 的固定斜面上，物体在水平推力 $F=50\UN$ 的作用下沿斜面向上运动，物体与斜面间的摩擦系数为 $\mu=0.1$，求物体运动的加速度（$g$ 取 $10\Umsq$）。
\onepicture[w=6.5cm, align=right]{\includesvg{figs/26}}

%% answer: $a$ = 2.3 m/s$^{2}$，方向沿斜面向上
\jdanswer{
$a=2.3\Umsq$，方向沿斜面向上
}
%% memo:
\memoanswer{
（1）画出示力图：示力图中缺一力不给分。

\onepicture[h=4.5cm]{\includesvg{figs/26_an}}

（说明：各力方向必须大致正确，力的大小比例不作要求）

（2）$f=\mu(mg\cos\alpha+F\sin\alpha)$，2 分。

（3）$F\cos\alpha-f-mg\sin\alpha=ma$，3 分。

$a=2.3\Umsq$，方向沿斜面向上（或在图上标出），方向、数值、单位各给 1 分。
}

\item
%% number: 27
%% paperName: 1989年上海高考第 27 题 14 分
%% typeId: 6
%% type: 计算题
%% chapter: 电磁感应
%% point: 电磁感应电路分析
%% method: 无
%% score: 14
%% degree: 450
%% duplicateId: 0
%% body:
在一磁感应强度 $B=0.5\UT$ 的匀强磁场中，垂直于磁场方向水平放置着两根相距为 $h=0.1\Um$ 的平行金属导轨 MN 与 PQ，导轨的电阻忽略不计。在两根导轨的端点 N、Q 之间连接一阻值 $R=0.3\UO$ 的电阻。导轨上跨放着一根长为 $L=0.2\Um$，每米长电阻 $r=2.0\UOpm$ 的金属棒 ab。金属棒与导轨正交放置，交点为 c、d。当金属棒以速度 $v=4.0\Ums$ 向左作匀速运动时，试求：
\begin{enumerate}
	\item
	电阻 $R$ 中的电流强度大小和方向；
	\item
	使金属棒作匀速运动的外力；
	\item
	金属棒 ab 两端点间的电势差。
\end{enumerate}
\onepicture[w=8cm, align=right]{\includesvg{figs/27}}

%% answer: （1）I ＝ 0.4 A，I 流向 N→Q；（2）F ＝ 0.02 N，方向向左；（3）U_ab ＝ 0.32 V
\jdanswer{
\begin{enumerate}
	\item
	$I=0.4\UA$，$I$ 流向 N$\to$Q
	\item
	$F=0.02\UN$，方向向左。
	\item
	$U_{ab}=0.32\UV$
\end{enumerate}
}
%% memo:
\memoanswer{
（1）$E_{cd}=hvB=0.1\times4.0\times0.5\UV=0.2\UV$，2 分。

$I=\frac{E_{cd}}{R+rh}=0.4\UA$，2 分。

$I$ 流向 N$\to$Q（或在图中标出），1 分。

（2）$F=IhB=0.4\times0.1\times0.5\UN=0.02\UN$，2 分。

外力 $F'=-F$，大小为 $0.02\UN$，方向向左。外力大小、方向各给 1 分。

（3）$E_{ab}=BLv=0.2\times4\times0.5\UV=0.4\UV$，给 1 分。

$U_{a}-U_{b}=E_{ab}-Irh=0.32\UV$，给 4 分。
}

\item
%% number: 28
%% paperName: 1989年上海高考第 28 题 15 分
%% typeId: 6
%% type: 计算题
%% chapter: 运动和力
%% point: 动量守恒定律
%% method: 分情况讨论
%% score: 15
%% degree: 450
%% duplicateId: 0
%% body:
有两块大小不同的圆形薄板（厚度不计），质量分别为 $M$ 和 $m$，半径分别为 $R$ 和 $r$，两板之间用一根长为 $0.4\Um$ 的轻绳相连结。开始时，两板水平放置并叠合在一起，静止于高度为 $0.2\Um$ 处（如图 \ref{1989上海28a}）。然后自由下落到一固定支架 C 上，支架上有一半径为 $R'$（$r<R'<R$）的圆孔，圆孔与两薄板中心均在圆孔中心轴线上，大板与支架发生没有机械能损失的碰撞。碰撞后，两板即分离，直到轻绳绷紧。在轻绳绷紧瞬间，两物体具有共同速度 $V$（如图 \ref{1989上海28b}）。问：
\begin{enumerate}
	\item
	若 $M=m$，则 $V$ 值为多大？
	\item
	若 $\frac{M}{m}=k$，试讨论 $V$ 的方向与 $k$ 值间关系。
\end{enumerate}
\twopicture[numA=1, numB=2, labelA=1989上海28a, labelB=1989上海28b, w=5.5cm, align=right]{\includesvg{figs/28a}}{\includesvg{figs/28b}}

%% answer: （1）V = 1 m/s；（2）k < 3 时 V > 0 两板向下运动；k > 3 时 V < 0 两板向上运动；k = 3 时 V = 0 两板瞬时静止
\jdanswer{
\begin{enumerate}
	\item
	$V=1\Ums$
	\item
	当 $k<3$ 时，$V>0$，两板向下运动；当 $k>3$ 时，$V<0$，两板向上运动；当 $k=3$ 时，$V=0$，两板瞬时静止。
\end{enumerate}
}
%% memo:
\memoanswer{
（1）$M$、$m$ 与固定支架碰撞前速度 $v_{0}=2\Ums$，给 2 分。

（2）碰撞后，$M$ 作初速 $v_{0}$ 向上的匀减速运动，$m$ 作初速 $v_{0}$ 向下的匀加速运动，$v_{1}=v_{0}-gt_{1}$；$v_{2}=v_{0}+gt_{2}$。

绳绷直时，$t_{1}=t_{2}$；$v_{1}+v_{2}=2v_{0}=4\Ums$（1），给 1 分。

$v_{1}^{2}=v_{0}^{2}-2gs_{1}$；$v_{2}^{2}=v_{0}^{2}+2gs_{2}$。

绳绷直时，$s_{1}+s_{2}=0.4\Um$；$v_{2}^{2}-v_{1}^{2}=2g(s_{1}+s_{2})=8\Um^{2}/\Us^{2}$（2），给 1 分。

由（1）（2）式解得：$v_{2}=3\Ums$；$v_{1}=1\Ums$，各给 1 分。

（说明：用其他方法求 $v_{1}$、$v_{2}$ 出正确值各给 2 分）

（3）在绳绷直过程中时间极短，重力的冲量忽略不计。

$mv_{2}-Mv_{1}=(m+M)V$，给 2 分。

①当 $M=m$ 时；$V=\frac{v_{2}-v_{1}}{2}=1\Ums$，给 1 分。

②当 $\frac{M}{m}=k$ 时；$V=\frac{v_{2}-\frac{M}{m}v_{1}}{1+\frac{M}{m}}=\frac{3-k}{1+k}$，给 2 分。

讨论：当 $k<3$ 时；$V>0$ 两板向下运动，给 2 分。

$k>3$ 时；$V<0$ 两板向上运动，给 1 分。

$k=3$ 时；$V=0$ 两板瞬时静止，给 1 分。

（说明：讨论得出一个结论得 2 分；讨论得出二个结论给 3 分）
}

\end{enumerate}

\end{document}
